\PassOptionsToPackage{table}{xcolor}
\documentclass[10pt,aspectratio=169]{beamer}
\newcommand{\LectureNumber}{28}
\input{course-theme.tex}
\title{Files and streams}
\subtitle{CSC103 Programming Fundamentals / Lecture 28}
\author{Dr. Muhammad Farhan}
\institute{Associate Professor of Computer Science\\Department of Computer Science\\COMSATS University Islamabad\\Sahiwal Campus}
\date{Fall 2026}
\begin{document}
\begin{frame}\titlepage\end{frame}

\begin{frame}[fragile,t]{The problem for this lecture}
\begin{columns}[T,onlytextwidth]\begin{column}{0.60\textwidth}
Create a File object for data/marks.txt. Display its absolute path, existence and whether it is a regular file. Explain environment-dependent output.
\end{column}\begin{column}{0.35\textwidth}\small
Represent data/marks.txt using a File object.
\end{column}\end{columns}
\note{Introduce the target problem before explaining the tools. Revisit it in the guided solution later.\par Syllabus reading: Liang Ch 12. CLO-3.}
\end{frame}

\begin{frame}[fragile,t]{Learning objectives}
\begin{columns}[T,onlytextwidth]\begin{column}{0.60\textwidth}
\begin{itemize}
\item Distinguish files, paths and streams
\item Inspect file metadata
\item Resolve relative paths against a working directory
\end{itemize}
\end{column}\begin{column}{0.35\textwidth}\small
CLO-3

Files and paths\par The File class\par Byte and character streams\par File operations and failures
\end{column}\end{columns}
\note{Explain how each objective supports the opening problem. The final questions check these objectives.\par Syllabus reading: Liang Ch 12. CLO-3.}
\end{frame}

\begin{frame}[fragile,t]{Files and paths}
\begin{itemize}
\item A file stores persistent data outside the running program.
\item A path identifies a location in a filesystem.
\item Relative paths depend on the process working directory.
\end{itemize}
\vspace{1mm}\begin{block}{Example}
\begin{lstlisting}
marks.txt
data/marks.txt

Both are relative paths.
\end{lstlisting}
\end{block}
\note{A program launched from an IDE can use a different working directory from the terminal. A filename alone does not identify a universal location.\par Syllabus reading: Liang Ch 12. CLO-3.}
\end{frame}

\begin{frame}[fragile,t]{The File class}
\begin{itemize}
\item java.io.File represents a pathname.
\item Constructing a File object does not create the physical file.
\item Methods such as exists and isFile inspect current metadata.
\end{itemize}
\vspace{1mm}\begin{block}{Example}
\begin{lstlisting}
java.io.File f = new java.io.File("marks.txt");
boolean present = f.exists();
\end{lstlisting}
\end{block}
\note{Metadata can change after checking it. The later I/O operation can still fail, so retain exception handling.\par Syllabus reading: Liang Ch 12. CLO-3.}
\end{frame}

\begin{frame}[fragile,t]{Byte and character streams}
\begin{itemize}
\item Byte streams move bytes. Character streams decode or encode text.
\item Text encoding connects stored bytes to characters.
\item Opening a stream acquires a resource that needs closing.
\end{itemize}
\vspace{1mm}\begin{block}{Example}
\begin{lstlisting}
Text: marks and names in UTF-8
Binary: an image or audio file

Use APIs appropriate to the data.
\end{lstlisting}
\end{block}
\note{Scanner parses text tokens. PrintWriter formats text. Avoid presenting all files as text readable by a character scanner.\par Syllabus reading: Liang Ch 12. CLO-3.}
\end{frame}

\begin{frame}[fragile,t]{File operations and failures}
\begin{itemize}
\item A missing file, a directory path and insufficient permission are different conditions.
\item Reading and writing have different effects.
\item Opening output can overwrite existing content.
\end{itemize}
\vspace{1mm}\begin{block}{Example}
\begin{lstlisting}
Before writing:
Choose a dedicated output path.
Decide overwrite or append behaviour.
Handle the opening failure.
\end{lstlisting}
\end{block}
\note{These examples use dedicated demo filenames. They do not enumerate or alter unrelated user files. Filesystem permissions remain external to the Java program.\par Syllabus reading: Liang Ch 12. CLO-3.}
\end{frame}

\begin{frame}[fragile,t]{An absolute path explains where Java looks}
\begin{itemize}
\item A relative pathname resolves against the process working directory.
\item Printing an absolute form is a useful diagnosis for a missing-file report.
\item The same source program can resolve a relative name differently under another launcher.
\end{itemize}
\vspace{1mm}\begin{block}{Example}
\begin{lstlisting}
Relative: data/marks.txt
Working directory: a course practice folder
Resolved location: that folder plus data/marks.txt
\end{lstlisting}
\end{block}
\note{Explain the mechanism using the displayed example. A relative pathname resolves against the process working directory. Printing an absolute form is a useful diagnosis for a missing-file report. The same source program can resolve a relative name differently under another launcher.\par Syllabus reading: Liang Ch 12. CLO-3.}
\end{frame}

\begin{frame}[fragile,t]{Metadata is an observation at one moment}
\begin{itemize}
\item exists reports whether the path exists when the call checks it.
\item isFile distinguishes an ordinary file from a directory.
\item A later operation still needs error handling because state or permissions can change.
\end{itemize}
\vspace{1mm}\begin{block}{Example}
\begin{lstlisting}
exists=false does not create a file.
exists=true does not guarantee that a later read succeeds.
\end{lstlisting}
\end{block}
\note{Explain the mechanism using the displayed example. exists reports whether the path exists when the call checks it. isFile distinguishes an ordinary file from a directory. A later operation still needs error handling because state or permissions can change.\par Syllabus reading: Liang Ch 12. CLO-3.}
\end{frame}


\begin{frame}[fragile,t]{A path is interpreted from a working directory}
\begin{center}\begin{tikzpicture}
\node[box] (n0) at (0.0,0) {Working directory};
\node[box] (n1) at (3.45,0) {Relative path};
\node[box] (n2) at (6.9,0) {Resolved location};
\node[alt] (n3) at (10.350000000000001,0) {File metadata};
\draw[arr] (n0) -- (n1);
\draw[arr] (n1) -- (n2);
\draw[arr] (n2) -- (n3);
\end{tikzpicture}\end{center}
\begin{block}{What to notice}Constructing a File describes a path; it does not create the referenced file.\end{block}
\note{Trace each labelled connection. Constructing a File describes a path; it does not create the referenced file.}
\end{frame}
\begin{frame}[fragile,t]{Key distinctions}
\small\renewcommand{\arraystretch}{1.5}
\rowcolors{2}{pale}{white}
\begin{tabularx}{\textwidth}{>{\raggedright\arraybackslash}X>{\raggedright\arraybackslash}X>{\raggedright\arraybackslash}X}
\rowcolor{navy}
\color{white}\bfseries File method & \color{white}\bfseries Question & \color{white}\bfseries Result type\\
getName() & Final pathname component? & String\\
getAbsolutePath() & Where does this path resolve? & String\\
exists() & Does the path currently exist? & boolean\\
isFile() & Does it identify a regular file? & boolean\\
\end{tabularx}
\vspace{3mm}\footnotesize These distinctions determine which operation or control structure fits the problem.
\note{\par Syllabus reading: Liang Ch 12. CLO-3.}
\end{frame}

\begin{frame}[fragile,t]{First example: execution}
\begin{lstlisting}
java.io.File file = new java.io.File("marks.txt");
System.out.println(file.getName());
System.out.println(file.isAbsolute());
\end{lstlisting}
\note{This is the main body from Java\_Examples/Lecture28.java. The source file includes the complete class. Predict the result before the next slide.\par Syllabus reading: Liang Ch 12. CLO-3.}
\end{frame}

\begin{frame}[fragile,t]{First example: result}
\begin{columns}[T,onlytextwidth]\begin{column}{0.60\textwidth}
marks.txt\par false\par No physical file is created by the constructor.\par The relative path will resolve in the program working directory.
\end{column}\begin{column}{0.35\textwidth}\small
Byte and character streams

File operations and failures
\end{column}\end{columns}
\note{Expected output and interpretation: marks.txt
false
No physical file is created by the constructor.
The relative path will resolve in the program working directory.\par Syllabus reading: Liang Ch 12. CLO-3.}
\end{frame}

\begin{frame}[fragile,t]{Guided practice}
\begin{columns}[T,onlytextwidth]\begin{column}{0.60\textwidth}
Create a File object for data/marks.txt. Display its absolute path, existence and whether it is a regular file. Explain environment-dependent output.
\end{column}\begin{column}{0.35\textwidth}\small
Attempt the solution before continuing.

Use the definitions and examples from this lecture.
\end{column}\end{columns}
\note{Give students 8 minutes to attempt the problem. The following slides provide a complete solution. Original solution guidance: Use getAbsolutePath(), exists() and isFile(). The exact path and Boolean results depend on where the program runs and what exists there.\par Syllabus reading: Liang Ch 12. CLO-3.}
\end{frame}

\begin{frame}[fragile,t]{Solution strategy}
\begin{columns}[T,onlytextwidth]\begin{column}{0.60\textwidth}
\begin{itemize}
\item Represent data/marks.txt using a File object.
\item Print the absolute path to expose working-directory resolution.
\item Inspect existence and file type without creating or modifying the path.
\end{itemize}
\end{column}\begin{column}{0.35\textwidth}\small
The next program uses fixed inputs so its result can be reproduced.
\end{column}\end{columns}
\note{Discuss why each step is needed. State the input assumptions before executing the program.\par Syllabus reading: Liang Ch 12. CLO-3.}
\end{frame}

\begin{frame}[fragile,t]{Complete solution}
\begin{lstlisting}
public class Solution28 {
  public static void main(String[] args) throws Exception {
    java.io.File file = new java.io.File("data/marks.txt");
    System.out.println(file.getAbsolutePath());
    System.out.println(file.exists());
    System.out.println(file.isFile());
  }
}
\end{lstlisting}
\note{Complete runnable file: Java\_Examples/Solution28.java. Inputs are fixed in the program.  \par Syllabus reading: Liang Ch 12. CLO-3.}
\end{frame}

\begin{frame}[fragile,t]{Execution trace}
\small\renewcommand{\arraystretch}{1.5}
\rowcolors{2}{pale}{white}
\begin{tabularx}{\textwidth}{>{\raggedright\arraybackslash}X>{\raggedright\arraybackslash}X>{\raggedright\arraybackslash}X}
\rowcolor{navy}
\color{white}\bfseries Observation & \color{white}\bfseries In an empty practice folder & \color{white}\bfseries Reason\\
Absolute path & Working folder + data/marks.txt & Relative path resolution\\
exists() & false & No file created\\
isFile() & false & Missing path is not a regular file\\
\end{tabularx}
\vspace{3mm}\footnotesize Follow the changing state in execution order.
\note{Manually derived trace for the complete solution. Represent data/marks.txt using a File object. Print the absolute path to expose working-directory resolution. Inspect existence and file type without creating or modifying the path.\par Syllabus reading: Liang Ch 12. CLO-3.}
\end{frame}

\begin{frame}[fragile,t]{Solution result}
\begin{columns}[T,onlytextwidth]\begin{column}{0.60\textwidth}
[working directory]\textbackslash{}data\textbackslash{}marks.txt\par false\par false
\end{column}\begin{column}{0.35\textwidth}\small
The absolute path and metadata depend on the working directory and existing files. This example assumes an empty practice folder.
\end{column}\end{columns}
\note{Expected result: [working directory]\textbackslash{}data\textbackslash{}marks.txt
false
false\par Syllabus reading: Liang Ch 12. CLO-3.}
\end{frame}

\begin{frame}[fragile,t]{A common incorrect approach}
new java.io.File("report.txt");\par // Student expects an empty report.txt to appear.
\vspace{1mm}\begin{block}{Example}
\begin{lstlisting}
The constructor only represents a pathname. A file-creation or writing operation is needed to create the physical file.
\end{lstlisting}
\end{block}
\note{Ask students to predict the failure before discussing the explanation. The right side gives the correction.\par Syllabus reading: Liang Ch 12. CLO-3.}
\end{frame}

\begin{frame}[fragile,t]{Boundary and failure checks}
\begin{columns}[T,onlytextwidth]\begin{column}{0.60\textwidth}
Try the program with a missing path, an existing text file and an existing directory.
\end{column}\begin{column}{0.35\textwidth}\small
Use getAbsolutePath(), exists() and isFile(). The exact path and Boolean results depend on where the program runs and what exists there.
\end{column}\end{columns}
\note{Use the specification to derive each expected result. Exercise solution: Use getAbsolutePath(), exists() and isFile(). The exact path and Boolean results depend on where the program runs and what exists there.\par Syllabus reading: Liang Ch 12. CLO-3.}
\end{frame}

\begin{frame}[fragile,t]{Check your understanding}
\begin{columns}[T,onlytextwidth]\begin{column}{0.60\textwidth}
Does a relative path start at the source-file folder automatically? What closes a stream reliably?
\end{column}\begin{column}{0.35\textwidth}\small
Write a prediction and explain the rule that supports it.
\end{column}\end{columns}
\note{Answer: No, it starts from the working directory. A try-with-resources statement closes an AutoCloseable resource on normal and exceptional exit.\par Syllabus reading: Liang Ch 12. CLO-3.}
\end{frame}

\begin{frame}[fragile,t]{Answer and explanation}
\begin{columns}[T,onlytextwidth]\begin{column}{0.60\textwidth}
No, it starts from the working directory. A try-with-resources statement closes an AutoCloseable resource on normal and exceptional exit.
\end{column}\begin{column}{0.35\textwidth}\small
The constructor only represents a pathname. A file-creation or writing operation is needed to create the physical file.
\end{column}\end{columns}
\note{Reconnect the answer to the relevant definition rather than treating it as a fact to memorise.\par Syllabus reading: Liang Ch 12. CLO-3.}
\end{frame}


\begin{frame}[fragile,t]{Compare cases and predict outcomes}
\small\renewcommand{\arraystretch}{1.5}\rowcolors{2}{pale}{white}
\begin{tabularx}{\textwidth}{>{\raggedright\arraybackslash}X>{\raggedright\arraybackslash}X>{\raggedright\arraybackslash}X}\rowcolor{navy}
\color{white}\bfseries Operation & \color{white}\bfseries Creates a file? & \color{white}\bfseries What it provides\\
new File("marks.txt") & No & Path object\\
getAbsolutePath() & No & Absolute path text\\
exists() & No & Existence observation\\
\end{tabularx}\par\vspace{3mm}\begin{exampleblock}{Discuss}Explain each expected result before running the example. Which case would reveal a wrong assumption?\end{exampleblock}
\note{Read the first two columns and invite predictions before discussing the outcome.}
\end{frame}

\begin{frame}[fragile,t]{Apply the idea}
\begin{alertblock}{Independent practice}Explain why new File("marks.txt") followed by exists() may be false. Show how to print the resolved path for diagnosis.\end{alertblock}\vspace{3mm}\begin{enumerate}\item Work through the input by hand.\item Write or correct the program.\item Justify the expected result.\end{enumerate}\note{Allow 3–5 minutes, then compare answers in pairs. Explain why new File("marks.txt") followed by exists() may be false. Show how to print the resolved path for diagnosis.}
\end{frame}

\begin{frame}[fragile,t]{Solution and reasoning}
\begin{lstlisting}
java.io.File file = new java.io.File("marks.txt");
System.out.println(file.getAbsolutePath());
System.out.println(file.exists());
\end{lstlisting}\vspace{1mm}\small\begin{itemize}
\item A path object does not imply that storage exists at that location.
\item The working directory can differ between a terminal and an IDE.
\item The printed path and Boolean depend on the actual run environment.
\end{itemize}\note{Answer: A path object does not imply that storage exists at that location. The working directory can differ between a terminal and an IDE. The printed path and Boolean depend on the actual run environment.}
\end{frame}

\begin{frame}[fragile,t]{Creating a file is different from naming it}
\begin{itemize}
\item Constructing File creates a path object; createNewFile attempts to create storage.
\item createNewFile returns false if the named file already exists.
\item Creation may throw IO\allowbreak Exception; an existence check does not guarantee later creation will succeed.
\end{itemize}
\vspace{3mm}\footnotesize\textcolor{accent}{Source: supplied ISB campus slides. See speaker notes and ISB\_Source\_Map.md.}
\note{Adapted from the supplied ISB campus folder: Java\_File\_Handling\_\_A\_Beginner\_s\_Guide.pptx, slides 4, 5; File Handling.pptx, slides 7, 10. Uses a fresh temporary directory for safe repeatable execution; the demonstration leaves that temporary file available for inspection.}
\end{frame}

\begin{frame}[fragile,t]{Worked problem: Creating a file is different from naming it}
\begin{block}{Task}Create a private temporary directory and call createNewFile twice on the same filename.\end{block}
\small Input: Values are fixed in the program for a reproducible demonstration.\par\vspace{3mm}\begin{exampleblock}{Before running}Predict the output and identify the variables that change.\end{exampleblock}
\note{Adapted from the supplied ISB campus folder: Java\_File\_Handling\_\_A\_Beginner\_s\_Guide.pptx, slides 4, 5; File Handling.pptx, slides 7, 10. Uses a fresh temporary directory for safe repeatable execution; the demonstration leaves that temporary file available for inspection.}
\end{frame}

\begin{frame}[fragile,t]{Complete ISB example (1/1)}
\begin{lstlisting}
public class IsbLecture28 {
  public static void main(String[] args) throws Exception {
    java.nio.file.Path folder = java.nio.file.Files.createTempDirectory("csc103-isb-");
    java.io.File file = folder.resolve("first.txt").toFile();
    System.out.println(file.createNewFile());
    System.out.println(file.createNewFile());
    System.out.println(file.isFile());
  }

}
\end{lstlisting}
\note{Adapted from the supplied ISB campus folder: Java\_File\_Handling\_\_A\_Beginner\_s\_Guide.pptx, slides 4, 5; File Handling.pptx, slides 7, 10. Uses a fresh temporary directory for safe repeatable execution; the demonstration leaves that temporary file available for inspection. Complete file: ISB\_Java\_Examples/IsbLecture28.java.}
\end{frame}

\begin{frame}[fragile,t]{Worked trace and expected result}
\small\renewcommand{\arraystretch}{1.4}\rowcolors{2}{pale}{white}
\begin{tabularx}{\textwidth}{>{\raggedright\arraybackslash}X>{\raggedright\arraybackslash}X>{\raggedright\arraybackslash}X}\rowcolor{navy}
\color{white}\bfseries Operation & \color{white}\bfseries Return / observation & \color{white}\bfseries Explanation\\
First createNewFile & true & New file created\\
Second createNewFile & false & File already exists\\
isFile & true & Path names a regular file\\
\end{tabularx}\par\vspace{2mm}
\begin{block}{Expected console output}\ttfamily\footnotesize true\par false\par true\end{block}
\note{Adapted from the supplied ISB campus folder: Java\_File\_Handling\_\_A\_Beginner\_s\_Guide.pptx, slides 4, 5; File Handling.pptx, slides 7, 10. Uses a fresh temporary directory for safe repeatable execution; the demonstration leaves that temporary file available for inspection. Trace and expected values independently worked for this edition.}
\end{frame}

\begin{frame}[fragile,t]{Extend the ISB example}
\begin{alertblock}{Practice}Why is an absolute path not universally wrong, despite the source warning against it?\end{alertblock}\vspace{3mm}\begin{exampleblock}{Explain your reasoning}Absolute paths can be appropriate when explicitly supplied or configured. Hard{-}coded machine{-}specific paths reduce portability.\end{exampleblock}
\note{Adapted from the supplied ISB campus folder: Java\_File\_Handling\_\_A\_Beginner\_s\_Guide.pptx, slides 4, 5; File Handling.pptx, slides 7, 10. Uses a fresh temporary directory for safe repeatable execution; the demonstration leaves that temporary file available for inspection. Answer: Absolute paths can be appropriate when explicitly supplied or configured. Hard{-}coded machine{-}specific paths reduce portability.}
\end{frame}
\begin{frame}[fragile,t]{Lecture summary}
\begin{columns}[T,onlytextwidth]\begin{column}{0.60\textwidth}
\begin{itemize}
\item files, paths and streams
\item Inspect file metadata
\item relative paths against a working directory
\end{itemize}
\end{column}\begin{column}{0.35\textwidth}\small
\begin{itemize}
\item Represent data/marks.txt using a File object.
\item Print the absolute path to expose working-directory resolution.
\item Inspect existence and file type without creating or modifying the path.
\end{itemize}
\end{column}\end{columns}
\note{Ask students to give one concrete example for the concepts listed. Use the worked solution to connect the topics.\par Syllabus reading: Liang Ch 12. CLO-3.}
\end{frame}

\begin{frame}[fragile,t]{After-class practice}
\begin{columns}[T,onlytextwidth]\begin{column}{0.60\textwidth}
Inspect a dedicated practice file and a directory using File methods. Record differences without deleting either.
\end{column}\begin{column}{0.35\textwidth}\small
Syllabus reading\par Liang Ch 12

Next lecture: Reading and writing text files
\end{column}\end{columns}
\note{Reading labels follow the supplied outline. Check topic headings against the available textbook edition. Require a program and expected results for the practice task.\par Syllabus reading: Liang Ch 12. CLO-3.}
\end{frame}
\end{document}
